\section{Sources, stopping scales, and correlation limits}
\label{sec:correlations}

The background flow has selected a temperature, but has not yet identified
any interacting correlation. We now insert finitely many sources into the
exact partition-function identity. Locality is important twice: it makes
each one-source normalization independent of the domain, and it prevents
distinct ideal sources from interacting before a suitably chosen stopping
scale. The remaining critical reference measure retains all of their
long-distance correlations.

\subsection{Finite source jets and scalar prefactors}

Let the source labels be $\{1,\ldots,q\}$, where $q\geq1$ is fixed.
The microscopic observable $X_\ell$ is either a spin at its marked site or
one of the four incident bond products, and its dimension is respectively
$d_\ell=1/8$ or $d_\ell=1$. Write
\[
 t_S=\prod_{\ell\in S}t_\ell,
 \qquad d(S)=\sum_{\ell\in S}d_\ell
 \qquad(\varnothing\ne S\subseteq\{1,\ldots,q\}).
\]
We calculate in the finite algebra
\[
 \mathcal A_q=\mathbb C[t_1,\ldots,t_q]/(t_1^2,\ldots,t_q^2).
\]
Thus coefficients are ordinary mixed derivatives at zero, with each label
used at most once. Products with overlapping label sets vanish. Exponentials
and logarithms with zero and unit constant term, respectively, are finite
polynomials in this algebra. In particular, no convergence assertion in
the source amplitudes is hidden in the use of a finite jet.

In a finite system $F$, let $H_j^F$ be its nonscalar background at scale
$r_j$, and let $\nu_j^F$ be the critical nearest-neighbor reference measure
tilted by its total representative. A source coefficient $U_j^S$ is a
marked list, with every carrier containing all the marks in $S$. We use
the same notation for the absolutely convergent sum of its functions.
Scalar coefficients are denoted by $c_j^S$. The exact identity at scale
$j$ is
\begin{equation}
 \begin{split}
 G_j^F(t)
 &:={\E}_{\lambda,F}
       \exp\!\left(\sum_{\ell=1}^q t_\ell r_j^{d_\ell}X_\ell\right)\\
 &=\exp\!\left(\sum_{\varnothing\ne S}c_j^S t_S\right)
   {\E}_{\nu_j^F}
       \exp\!\left(\sum_{\varnothing\ne S}U_j^S t_S\right).
 \end{split}
 \label{src:partition-jet}
\end{equation}
Here the expectation on the first line is at the critical branch from
\Cref{flow:branch}, with the finite-system prescription under consideration.
The preliminary transformation establishes this identity at $j=0$.

For precision, apply the local map to
$H_j^F+\sum_S U_j^S t_S$. If its newly produced nonscalar and scalar
source coefficients are $\widehat U_j^S$ and $a_j^S$, then
\begin{equation}
 U_{j+1}^S=L^{d(S)}\widehat U_j^S,
 \qquad
 c_{j+1}^S=L^{d(S)}(c_j^S+a_j^S).
 \label{src:source-recursion}
\end{equation}
These factors express the substitution
$t_\ell\mapsto L^{d_\ell}t_\ell$ in $G_j^F$.
They are compatible with products because $d(S)$ is additive on disjoint
label sets. The scalar prefactor is removed before the next local map is
applied. It is therefore only rescaled and added to; it never becomes an
input to a nonlinear coefficient. Scalar background contributions cancel
between numerator and denominator in the normalized expectation.

The marked conclusions of \Cref{rg:map} apply to these operations. For a
fixed $q$, the map on the source coefficients is polynomial, with bounded
derivatives on bounded sets of marked norms, uniformly in the scale as
long as the background stays in its fixed small ball. This follows either
by decorating the convergent list expansion with the labels, or by Cauchy
estimates on a sufficiently small source polydisc at that background.
The polydisc and the resulting constants may depend on $q$ and the source
bounds; the background ball does not. A one-source nonscalar map and its
scalar-increment functional differ from their zero-background versions
by at most $C_L\norm{H_j^F}_{r_j}$ in operator norm. Scalar increments
retain their generating carriers for estimates before those carriers are
discarded. In particular, their sums have no volume factor: any one of
their marks roots the entire list.

Let $\mathcal E$ be the set of energy labels. On logarithmic jets define
\[
 \mathcal Q K
   =K-\sum_{\ell\in\mathcal E}t_\ell[t_\ell]K.
\]
Then the exact centering identity is
\begin{equation}
 [t_1\cdots t_q]\exp(\mathcal Q\log G_j^F)
 =r_j^{d(\{1,\ldots,q\})}
   {\E}_{\lambda,F}
   \left[
    \prod_{\ell\notin\mathcal E}X_\ell
    \prod_{\ell\in\mathcal E}
       (X_\ell-{\E}_{\lambda,F}X_\ell)
   \right].
 \label{src:centering}
\end{equation}
Indeed, replacing an energy by its centered version subtracts precisely
$t_\ell r_j^{d_\ell}{\E}_{\lambda,F}X_\ell$ from the logarithm of its
generating function. By \eqref{src:partition-jet}, this coefficient is
\[
 [t_\ell]\log G_j^F
   =c_j^{\{\ell\}}+{\E}_{\nu_j^F}U_j^{\{\ell\}}
   =r_j^{d_\ell}{\E}_{\lambda,F}X_\ell.
\]
Thus centering removes both the extracted scalar and the residual
one-source mean. The scalar coefficients $c_j^{\{\ell\}}$ for energy
labels may be omitted throughout the estimates; applying $\mathcal Q$
to the remaining logarithm still removes the residual mean. All spin
singleton scalars, and all multiple-label scalars, must still be kept.
The actual domain and boundary-state means in
\eqref{src:centering} are not replaced by plane means.

\subsection{The common one-source normalizations}

Let $u_j^\ell$ be the nonscalar plane singleton list for the microscopic
source $X_\ell$ on the critical background trajectory $h_j$. Start with
the degree-one coefficient of the same preliminary transformation as
in \eqref{src:partition-jet}, normalized by $r_0^{d_\ell}$, and set
\[
 u_{j+1}^\ell=L^{d_\ell}D\mathcal R_j(h_j)u_j^\ell.
\]
Scalar terms are recorded separately as above; energy lists are taken
modulo constants. This definition uses the local list map, before any
interacting thermodynamic state is constructed.

\begin{proposition}[One-source limits]
\label{src:scalar}
There are real analytic functions $Z_\sigma(\lambda)$ and
$Z_\epsilon(\lambda)$ on a neighborhood of zero, with
\[
 Z_\sigma(0)=Z_\epsilon(0)=1,
\]
such that the plane one-source lists on the critical trajectory satisfy
\begin{equation}
 u_j^\ell
   =Z_\ell(\lambda)v_j^{(d_\ell)}+\eta_j^\ell,
 \qquad
 \norm{\eta_j^\ell}_{r_j,\bullet}\longrightarrow0.
 \label{src:one-source}
\end{equation}
Here $Z_\ell$ is $Z_\sigma$ for a spin and $Z_\epsilon$ for any incident
bond orientation; energy lists are understood modulo constants.
The convergence is locally uniform in the coupling. The two functions
are nonzero on one smaller coupling interval, independently of domains,
boundary conditions, and the number of sources.
\end{proposition}

\begin{proof}
Use the bounded splittings from \Cref{rg:linear}, with the distinguished
vectors from \Cref{rg:fields}. For a normalized one-source coefficient the
zero-background matrix is
\[
 A_j=
 \begin{pmatrix}1&B_j\\0&D_j\end{pmatrix},
 \qquad \norm{B_j}\leq C_L,
 \qquad \norm{D_j}\leq\rho<1.
\]
The derivative on the critical plane trajectory differs from $A_j$ by an
operator $E_j$ with
\[
 \norm{E_j}\leq C_L\norm{h_j}_{r_j}
                  \leq C|\lambda|\gamma^j.
\]
Products of consecutive $A_j$ are uniformly bounded: their lower-right
blocks contract geometrically and their upper-right blocks are bounded
by $C_L\sum_{s\geq0}\rho^s$. Expanding a perturbed product by the locations
of its factors $E_j$ therefore bounds it by a constant times
$\prod_j(1+C\norm{E_j})$. This product is finite, uniformly for the
coupling in a sufficiently small complex disc. Thus the one-source
trajectory is bounded.

Write its coordinates as $\alpha_j,w_j$. The stable recurrence gives
\[
 \norm{w_j}
 \leq \rho^j\norm{w_0}
       +C|\lambda|\sum_{i<j}\rho^{j-1-i}\gamma^i.
\]
This tends to zero and is summable in $j$; it is bounded by $C\theta^j$
for any fixed $\max(\rho,\gamma)<\theta<1$, after enlarging $C$.
Moreover
\[
 |\alpha_{j+1}-\alpha_j|
 \leq C_L\norm{w_j}+C|\lambda|\gamma^j.
\]
Hence $\alpha_j$ converges, and subtracting its limit times
$v_j^{(d_\ell)}$ proves \eqref{src:one-source}. The preliminary source
map and every later derivative are analytic in the coupling. The bounds
just given are uniform on a smaller complex disc, so the limiting scalar
is analytic there. Reality of all maps on real inputs gives real values
on the real axis. Spin-flip symmetry makes every plane spin singleton
scalar contribution identically zero.

Translation covariance makes the limit independent of the marked site.
Choose the energy splitting equivariantly under the square symmetries.
Rotation about the mark carries the four incident bond sources into one
another and fixes the scalar energy line, so their limit scalars agree.
At zero coupling the spin source is exactly on its distinguished line
after the preliminary step. The average of the four bond sources is
$e(z)$, so this average likewise starts exactly on the energy line.
The normalized line multiplier is one. Linearity of a one-source
derivative and equality of the four scalars give both asserted
intercepts. Continuity now gives simultaneous nonvanishing after one
decrease of the coupling interval. All choices were made on the plane
background and for the two field types, without reference to a domain
or a collection of other marks.
\end{proof}

\subsection{Separated sources before they interact}

We next consider the sites associated with a fixed configuration of
distinct macroscopic interior points. Let $\Delta_*$ be a lattice-unit
separation no larger than any pairwise mark distance or any mark-to-edge
distance. In the plane comparison on a torus, it is also much smaller
than the side $N$. It can be chosen comparable to $a^{-1}$ for each fixed
configuration. Bounded rounding errors at the marks are harmless after
decreasing this separation by a fixed amount.

Choose the last scale $r_k$ below a sufficiently small fixed multiple of
$\Delta_*$ and below the physical boundary-chart thresholds in
\Cref{flow:boundary}. On a torus also require $r_k$ below the chart
threshold in \Cref{flow:torus}. For a fixed domain and configuration,
and in the torus regime $N/\Delta_*\to\infty$, these choices give
\begin{equation}
 k\longrightarrow\infty,
 \qquad 0<c_1\leq a r_k\leq c_2<\infty.
 \label{src:stop-scale}
\end{equation}
The constants may depend on the fixed configuration and domain. The
small stopping fraction is chosen so that all ideal source supports and
their protective neighborhoods are disjoint and interior at every step
ending no later than $k$.

\begin{proposition}[Separated-source comparison]
\label{src:separated}
At the stopping scale just described, the exact finite-system source
coefficients satisfy
\begin{align}
 U_k^{\{\ell\}}
   &=Z_\ell v_k^{(d_\ell)}+o(1)
       &&\text{in marked norm}, \label{src:single-stop}\\
 \norm{U_k^S}_{r_k,\bullet}&=o(1)
       &&(|S|\geq2), \label{src:multiple-stop}\\
 c_k^S&=o(1)
       &&\bigl(|S|\geq2\text{ or }S
                         \text{ is a spin singleton}\bigr).
       \label{src:scalar-stop}
\end{align}
The ideal functions in \eqref{src:single-stop} are restricted to the
system, with their cutting squares contained in its interior. On tori
the conclusions hold as $a\to0$ and $N/\Delta_*\to\infty$.
All errors use the running normalization $r_k^{d(S)}$.
\end{proposition}

\begin{proof}
Let $\Pi_F$ denote restriction to intact carriers in a domain, or the
small-diameter periodic projection used in \Cref{flow:torus}. At each
scale use the trial coefficients
\[
 W_j^{\{\ell\}}=\Pi_F u_j^\ell,
 \qquad W_j^S=0\quad(|S|\geq2),
 \qquad \widetilde c_j^S=0
\]
for the scalar labels retained in \eqref{src:scalar-stop}. Let $\mathcal I$
be this collection of scalar label sets: all sets with at least two
labels, and every spin singleton. On source states use the finite product
norm
\[
 \|(U,c)\|_{j,\mathrm{src}}
   :=\sum_{\varnothing\ne S}\|U^S\|_{r_j,\bullet}
          +\sum_{S\in\mathcal I}|c^S|.
\]
Write $\mathsf U_j=(U_j,c_j)$ for the exact state and
$\mathsf W_j=(W_j,\widetilde c_j)$ for the trial state. Let $\Phi_j^F$
be the source step \eqref{src:source-recursion} with the actual
finite-system background $H_j^F$ held fixed. Thus no background error
is included among these source variables. Define the error and the
one-step trial defect by
\[
 \mathfrak e_j=\mathsf U_j-\mathsf W_j,
 \qquad
 \mathfrak d_j=\Phi_j^F(\mathsf W_j)-\mathsf W_{j+1}.
\]
Energy singleton scalars are omitted from these states, consistently
with \eqref{src:centering}.

The defect is bounded by
\begin{equation}
 \|\mathfrak d_j\|_{j+1,\mathrm{src}}
 \leq\varepsilon_j\exp(-c\Delta_*/r_{j+1}),
 \qquad
 \sup_j\varepsilon_j<\infty,
 \qquad \varepsilon_j\longrightarrow0.
 \label{src:defect}
\end{equation}
We verify both factors; a fixed small stopping fraction alone would not
make the exponential in \eqref{src:defect} tend to zero.

First, the list/log map has connected generating carriers. A coefficient
joining two distinct marks must span their separation. A discrepancy
between an interior singleton calculation and its plane calculation
must reach the physical edge; a term containing a boundary-family input
likewise connects its mark to the edge. In the torus comparison a
projection exception must have a nonlocal or winding generating
neighborhood. These are precisely diameter conditions covered by the
doubled-weight and marked estimates in \Cref{rg:map}. They give the
exponential factor in \eqref{src:defect}, including before scalar
projection. They also cover long unchanged singleton carriers.

For the vanishing factor, terms containing a background input have the
extra bound $C|\lambda|\gamma^j$, by \Cref{flow:branch,flow:boundary}.
The corresponding torus bounds follow from \Cref{flow:torus}; in
particular its background discrepancy has the stronger norm bound
\[
 C|\lambda|\gamma^j\exp(-c'N/r_j),
\]
which supplies a spatial gain even when that discrepancy is used as one
input without reestablishing its generating history.

For terms with no background input, decompose every trial singleton by
\Cref{src:scalar} into its ideal term and remainder. A term containing at
least one remainder has a factor tending to zero in marked norm. This
factor is retained simultaneously with the spatial exponential: the
latter is paid from the reserved carrier weights, whereas the former is
one of the multilinear input norms. A term containing only ideal
singletons cannot be exceptional or join distinct marks. Each such
singleton and its protective neighborhood has radius $O(r_{j+1})$
about its own mark. The stopping fraction makes these neighborhoods
pairwise disjoint and interior, so the logarithm of their factorizing
list weight has no multiple-label coefficient. Their singleton
transfers are exactly the plane ones.

There is no missing spin scalar in this comparison. In the plane it is
zero by parity. On intact interior carriers the local graph and the
available-spin representative projection are the same in the finite
system and on the plane. An interior ideal spin therefore produces no
scalar discrepancy. Every remaining spin singleton discrepancy reaches
the edge or involves the torus discrepancy or a remainder, and is
covered by the preceding estimates. Multiple-label scalar coefficients
have the same connected generating carriers as their nonscalar
counterparts, so the same proof applies to them. Energy singleton
scalars are deliberately absent from the error state by
\eqref{src:centering}. For example, after decreasing $c$ one may take a
bounded multiple of
$|\lambda|\gamma^j+\sum_\ell\norm{\eta_j^\ell}_{r_j,\bullet}$
as $\varepsilon_j$. This proves \eqref{src:defect}.

The preliminary defect obeys
$C\exp(-c\Delta_*/r_0)$, with a fixed prefactor, by the preliminary
local-list estimates. The microscopic sources are interior and
separated. On a unit neighborhood of the bounded trial jets, the source
maps have a scale-independent Lipschitz constant $K\geq1$. It includes
all the finite multipliers $L^{d(S)}$ for the retained scalars.
Consequently $\|\mathfrak e_l\|_{l,\mathrm{src}}$ at any $l\leq k$,
as long as the state stays in this neighborhood, is bounded by
\begin{equation}
 C K^l e^{-c\Delta_*/r_0}
 +C\sum_{i<l}K^{l-i-1}\varepsilon_i
                       e^{-c\Delta_*/r_{i+1}}.
 \label{src:error-sum}
\end{equation}
Put $q_* =\inf \Delta_*/r_k>0$. At $l=k$, reindexing the sum by
$h=k-i-1$ bounds it by
\[
 C\sum_{h=0}^{k-1}K^h\varepsilon_{k-h-1}e^{-c q_*L^h}.
\]
The majorant $C\sup_i\varepsilon_i\,K^h e^{-c q_*L^h}$ is summable.
For each fixed $h$ the corresponding $\varepsilon_{k-h-1}$ tends to
zero, so dominated convergence gives an $o(1)$ bound. The same bound is
uniform over $l\leq k$: for $l$ large use the same tail majorant and the
smallness of the late $\varepsilon_i$, while the finitely many early
$l$ have $\Delta_*/r_l\to\infty$. The preliminary term also tends to
zero uniformly, since geometric growth of $\Delta_*/r_0$ dominates
$K^k$. Thus the estimates bootstrap inside the unit neighborhood.
Together with \eqref{src:one-source}, they prove all three conclusions.
No rate of convergence for that remainder was needed.
\end{proof}

\subsection{Evaluation in a bounded domain}

At the terminal scale, the entire nonscalar background in a fixed
bounded domain has sup norm tending to zero. To see this directly,
cover the available lattice sites by $O(1)$ squares of radius $r_k$;
the number is bounded by \eqref{src:stop-scale}. Every surviving
nonconstant representative depends on an available spin, hence its
carrier meets one of these squares. Therefore
\begin{equation}
 \left\|\sum_A H_{k,A}^F\right\|_\infty
 \leq C_F\norm{H_k^F}_{r_k}
 \leq C_F|\lambda|\gamma^k\longrightarrow0.
 \label{src:total-background}
\end{equation}
The sum here includes both restricted bulk and boundary lists. The
covering argument remains valid for carriers extending outside the
domain, because a nonconstant term must still contain an available
variable spin.

All terminal backgrounds are real for real $\lambda$. Their tilts are
thus positive probability measures. If a real exponent has sup norm
at most $b$, its normalized density lies between $e^{-2b}$ and $e^{2b}$.
Equation \eqref{src:total-background} consequently permits replacing
$\nu_k^F$ by the ordinary critical reference measure in expectations
of uniformly bounded terminal tests, with error $o(1)$.

Taking the logarithm of the last factor in
\eqref{src:partition-jet} gives a finite polynomial in moments of the
source coefficients, through the degrees in use. Their marked norms
bound the sums of sup norms. \Cref{src:separated} and the preceding
probability comparison therefore give, coefficientwise,
\begin{equation}
 \mathcal Q\log G_k^F
 =\mathcal Q\log {\E}_{0,F}
      \exp\!\left(\sum_{\ell=1}^q
                t_\ell Z_\ell v_k^{(d_\ell)}\right)+o(1).
 \label{src:terminal-log}
\end{equation}
Here and below coefficientwise $o(1)$ refers to the finitely many
coefficients of positive degree. The scalar conclusions of
\Cref{src:separated} are essential in this identity; energy singleton
scalars disappear under $\mathcal Q$, while spin singleton scalars do
not disappear by convention.

The cutting squares of the ideal fields are disjoint and interior.
Condition on their exterior, retaining their cutting contours. The
critical nearest-neighbor Markov property makes their interiors
conditionally independent. For every subset of labels, insertion of
the conditional expectations $P_{r_k,z_\ell}\mathcal O_{d_\ell}(z_\ell)$ is
therefore exactly equivalent to insertion of the microscopic
$\mathcal O_{d_\ell}(z_\ell)$. This is an equality of all the subset moments,
not only of the full product. The representatives used for energy may
differ by constants, which leave $\mathcal Q\log$ unchanged.
Exponentiating \eqref{src:terminal-log} and using
\eqref{src:centering} identifies the requested centered product, at
running normalization, with the corresponding critical product of
spins and symmetric incident-bond averages, multiplied by
$\prod_\ell Z_\ell$, up to $o(1)$.

Passing to the normalization $a^{-d_\ell}$ multiplies the error by
$(a r_k)^{-d(\{1,\ldots,q\})}$, which is bounded by
\eqref{src:stop-scale}. The nearest-neighbor convergence in
\Cref{in:fields} applies to each term of the finite expansion of the
symmetric bond averages. Backward incident bonds amount to bounded
lattice shifts, and locally uniform convergence on separated bulk
configurations permits those shifts. Thus every orientation has the
same limiting scalar energy convention. This proves the bounded-domain
correlation assertion along the full mesh sequence; the quantity
$a r_k$ need not itself converge.

\subsection{The periodic thermodynamic limit}

The preceding total-background argument does not apply to tori whose
side is arbitrarily large compared with $r_k$. We instead use one
additional source without a scaling multiplier. Its contraction
controls local expectations even when the total perturbation is large.

\begin{samepage}
\begin{proposition}[Periodic limits and residual local laws]
\label{src:plane}
After one common decrease of the coupling interval, the following hold.
\begin{enumerate}
 \item At the selected temperature, every bounded observable of finitely
 many lattice spins has a limit along periodic square boxes as their
 side tends to infinity. In particular their finite-dimensional laws
 have a periodic thermodynamic limit at each fixed lattice spacing.
 \item Let $j\to\infty$ and $N/r_j\to\infty$. For each fixed $D<\infty$,
 the expectations of every test of sup norm at most one supported in a
 square of radius $D r_j$ under the residual tilted torus law
 $\nu_j^N$ and the critical nearest-neighbor torus law differ by
 $o(1)$, uniformly over these tests.
\end{enumerate}
\end{proposition}
\end{samepage}

\begin{proof}
We first establish an unscaled marked contraction. Split a short source
function into its spin-flip even and odd parts. Combining the response
and remainder bounds in \Cref{rad:operator} bounds the norm of its
conditional average modulo constants by $C(CM/L)^{1/8}$ times its
sup norm; the odd sector gives the slower decay. Here there is neither
a volume multiplier nor the multiplier $L^d$ used for a scaling field.
The doubled output
weight on its fixed-size carrier costs only a constant. Long source
carriers have the exponentially small tail from \Cref{rg:map}.
The total marked norm sums the list directly, so these estimates give
\begin{equation}
 \norm{D\mathcal R_j(0)u}_{r_{j+1},\bullet}
 \leq q_0\norm{u}_{r_j,\bullet},
 \qquad q_0<1,
 \label{src:unscaled-contraction}
\end{equation}
after including this condition when choosing the large block ratio.
Canonical subtraction of constants costs at most a factor two, already
included here. No parity restriction remains on $u$. The estimate
holds on tori through the chart scales in \Cref{flow:torus}, since
short cutting squares have the same local critical continuation and
long carriers are handled by their norm tail.

By \Cref{rg:map}, the derivative at a residual background differs by
at most $C_L\norm{H_i^N}_{r_i}$ from this derivative. The torus bound
in \Cref{flow:torus} and one decrease of $|\lambda|$ therefore give
a uniform contraction factor $q_1<1$ at every permitted step. Denote
the scalar-increment functional for this unscaled derivative by
$A_i(H_i^N)$. The same proposition gives
\begin{equation}
 \norm{A_i(H_i^N)}\leq C_L,
 \qquad
 \norm{A_i(H_i^N)-A_i(0)}
    \leq C_L\norm{H_i^N}_{r_i}.
 \label{src:unscaled-scalars}
\end{equation}

Suppose first that a bounded test $F$ is inserted at scale $j$ with
support in a square of radius $D r_j$. Subtract its uniform-spin mean
$c$ and use that square as its marked carrier. The initial marked norm
is at most a constant $C_D$. This carrier may initially be long; the
long-carrier part of \eqref{src:unscaled-contraction} applies to it.
Thus $D$ changes the initial bound, not the contraction factor or the
coupling interval. Propagate this unscaled derivative through
any terminal index $s$ with $r_s$ below the torus chart threshold.
Let $w_i$ be its nonscalar list at scale $i$. The exact first-derivative
version of the partition identity reads
\begin{equation}
 {\E}_{\nu_j^N}F
 =c+\sum_{i=j}^{s-1} A_i(H_i^N)w_i
      +{\E}_{\nu_s^N}w_s,
 \qquad
 \norm{w_i}_{r_i,\bullet}\leq C_D q_1^{i-j}.
 \label{src:local-telescope}
\end{equation}
The final expectation is bounded by its marked norm, using positivity.

Compare this calculation with the same one on the zero-background
trajectory, starting from the same $F$, and denote its lists by
$w_i^0$. The operator difference is bounded by $C|\lambda|\gamma^i$.
Using a common contraction factor $q_1$ in both trajectories gives
\[
 \norm{w_i-w_i^0}_{r_i,\bullet}
 \leq C_D|\lambda|\gamma^j(i-j)q_1^{i-j-1}
 \quad(i>j).
\]
Equations \eqref{src:unscaled-scalars} and
\eqref{src:local-telescope} can now be summed geometrically. The
difference of their scalar sums is at most $C_D|\lambda|\gamma^j$;
the two terminal expectations contribute at most $2C_Dq_1^{s-j}$.
Thus
\begin{equation}
 \left|{\E}_{\nu_j^N}F-{\E}_{0,N}F\right|
 \leq C_D\bigl(|\lambda|\gamma^j+q_1^{s-j}\bigr).
 \label{src:residual-comparison}
\end{equation}
Taking the last permitted $s$ makes $s-j\to\infty$ whenever
$N/r_j\to\infty$. This proves the second assertion, uniformly over
the stated tests.

For the first assertion, fix a bounded observable $F$ of finitely many
lattice sites, and choose a fixed connected carrier containing them.
Insert it at the microscopic scale. If this carrier does not fit a
preliminary cutting square, use the unchanged long-singleton rule,
with its protective neighborhood. Its fixed size only changes the
finite constants in the preliminary marked bound. Thereafter the
unscaled contraction and scalar estimate give a summable tail
$C_F\sum_{i\geq m}q_1^i$, uniformly in every sufficiently large torus
on which these scales are available.

For any fixed finite set of scales, the background coefficients
converge to their plane counterparts by \Cref{flow:torus}. The same
is true of the local derivative and each scalar increment. Indeed,
lists contained in a common chart agree with the plane list rule;
nonchart lists or omitted long inputs have a diameter comparable to
$N$ and vanish by the marked diameter estimates. The weighted tails
permit the passage from finite lists to the complete coefficient.
Thus the first $m$ scalar increments have a limit as $N\to\infty$.
The sum of subsequent increments and the final nonscalar expectation
is uniformly negligible as $m\to\infty$, by the contraction just
proved. Applying the telescope from the preliminary scale proves that
${\E}_{\lambda,N}F$ converges.

Taking $F$ to be cylinder indicators gives consistent limiting
finite-dimensional probabilities, hence the asserted thermodynamic
state. This construction does not assume uniqueness of other Gibbs
states or a choice of a symmetry-broken phase. Its coupling restriction
comes from a uniform one-mark operator bound and is independent of
$F$ and of the number of sites on which $F$ depends.
\end{proof}

\subsection{Completion of the universality theorem}

\begin{lemma}[Mixed reference correlations in the plane]
\label{src:reference-plane}
For every finite configuration of pairwise distinct physical insertion
points, use the spin and energy normalizations of
\eqref{model:spin}--\eqref{model:energy}, with each energy centered in its
actual nearest-neighbor state. The full-plane critical mixed correlations
then have full-sequence scaling limits. These limits are independent of
bond directions and equal the limits of their bounded-domain counterparts
along an exhaustion by fixed-plus disks, using the actual disk-state means
there.
\end{lemma}

\begin{proof}
Fix separated physical insertion points and a small positive stopping
radius below all their pairwise separations. Choose $k=k(a)$ so that
$a r_k$ stays between two positive constants below that radius.
At zero background, singleton transport is exactly nested conditional
averaging. Its nonscalar list is therefore
$u_k^\ell=r_k^{d_\ell}\pi P_{r_k,z_\ell}X_\ell$, also for an individual
directed bond. By \Cref{src:scalar,rg:fields} these lists have bounded
marked norm. Hence there is a constant $c_{a,\ell}$ such that
\[
 f_{a,\ell}=r_k^{d_\ell}P_{r_k,z_\ell}X_\ell-c_{a,\ell}
\]
has uniformly bounded sup norm.
The constants vanish for spins and are immaterial for centered energies.
The cutting squares are disjoint and their union lies in one fixed
physical square $B_B$. These functions are identical in the plane and
in every sufficiently large disk $D_T$ containing $B_B$: each conditional
kernel is determined inside its own cutting square.

The conditional Markov identity for every subproduct identifies the
original scaled centered correlation with
$(a r_k)^{-d(\{1,\ldots,q\})}$ times the corresponding centered product
of the $f_{a,\ell}$, in either system. Let $p_a$ and $p_{a,T}$ denote
these original nearest-neighbor correlations in the plane and in $D_T$
with plus boundary. Conditioning on a common surrounding square and
applying \Cref{in:ratio} with repeated critical RSW in the annuli between
$B_B$ and the distant wall gives a total-variation error $\omega(T)\to0$
on their common support, uniformly for sufficiently small $a$.
The uniform sup norms control all subset moments. Actual-state centering
is a finite polynomial in those moments, so
\[
 \limsup_{a\downarrow0}|p_a-p_{a,T}|\le C\omega(T).
\]
For every fixed $T$, bounded-domain convergence in \Cref{in:fields} gives
$p_{a,T}\to\mathcal C_{D_T}^{+}$. Consequently
\[
 \limsup_{a,b\downarrow0}|p_a-p_b|\le2C\omega(T).
\]
Sending $T$ to infinity proves the full mesh limit and identifies it with
$\lim_{T\to\infty}\mathcal C_{D_T}^{+}$. The same comparison proves
independence of the bond directions. The order is a fixed positive
stopping radius, then the mesh limit at each fixed $T$, and finally
$T\to\infty$. Only bounded-domain mixed convergence and the preceding
uniform conditional-field bounds have been used.
\end{proof}

\begin{proof}[Proof of \Cref{model:main}]
\label{src:main-proof}
Fix the allowed potential once and for all. \Cref{flow:branch} gives
the critical inverse-temperature branch with the specified intercept,
positive on a small real coupling interval. The reference perturbation
has real coefficients for either sign of the coupling; none of the
background estimates requires a sign condition on the finite-range
potential. \Cref{src:scalar} gives the common nonzero normalizations
with their unit intercepts. Choose the interval once so that the
background, boundary, torus, nonvanishing, and unscaled-contraction
bounds used above all hold. These choices do not involve a domain,
insertion count, or insertion configuration. Constants in a fixed
finite source jet may depend on those data, as the theorem permits.

For any fixed allowed bounded domain, boundary type, and separated
configuration, \Cref{src:separated} and the bounded-domain evaluation
prove the requested full-sequence limit. The microscopic energy
sources may have either specified direction. Their plane scalars
coincide, and the critical symmetric-average comparison has exactly
the common nearest-neighbor energy limit. Equation
\eqref{src:centering} centers every source in its actual finite-domain
state. For a fixed negative boundary, the same argument uses its
physical prescription throughout; no boundary-dependent normalization
has been introduced.

For the full-plane assertion, the first part of \Cref{src:plane}
establishes the required periodic-box limit at each fixed mesh.
Consider an arbitrary sequence $a\to0$ and choose its stopping scales
by \eqref{src:stop-scale}. On tori with $N/r_k\to\infty$,
\Cref{src:separated} still reduces the final source jet to the bounded
ideal singleton fields, with vanishing remaining scalar and
multiple-source coefficients. Products of any subset of these ideal
fields have uniformly bounded sup norms and support in a square of
radius $D r_k$, for a fixed $D$ depending on the configuration.
The second part of \Cref{src:plane} therefore replaces their residual
tilted expectations by critical torus expectations. This step does
not assert smallness of a total background on an expanding torus.

The latter local critical expectations can in turn be replaced by
plane expectations as $N/r_k\to\infty$. Surround their support by a
planar chart. The pointwise ratio estimate and repeated critical
RSW in the intervening annuli make the effect of arbitrary spins at
the distant chart wall tend to zero, uniformly for the bounded tests.
Both the torus and plane laws are mixtures of those continuations.
The cutting squares defining the ideal fields are identical local
critical squares in both systems. Their simultaneous conditional
expectation identity now recovers the original microscopic critical
spins and symmetric energies in the plane. \Cref{src:reference-plane}
finishes the same comparison there. The separately cited plane spin
formula identifies its nonzero two-spin amplitude.

To make the order of limits explicit, at each fixed mesh all unscaled
moments appearing in the centered product converge as $N\to\infty$,
by \Cref{src:plane}. The centered moment is a finite polynomial in
these moments. Choose $N$ large enough that its error after the fixed
mesh normalization is smaller than any prescribed sequence tending
to zero, and also so that $N/r_k\to\infty$. The joint comparison just
proved applies to these chosen boxes. Since the original mesh
sequence was arbitrary, this proves the stated iterated, full-sequence
plane limit. Its two-spin coefficient is nonzero: it is
$Z_\sigma(\lambda)^2$ times the positive nearest-neighbor plane
amplitude multiplying $|x-y|^{-1/4}$. In particular the selected
temperature branch has the massless power-law behavior required here;
that critical behavior was a conclusion, not an input to its
construction.

Finally let $\phi:D\to D'$ be as in the theorem. The nearest-neighbor
spin field and the energy field centered relative to the plane
background have their usual primary covariance, by \Cref{in:fields}.
If $m_D^b(y)$ denotes the domain one-point function of that energy
field, it has the same factor $|\phi'(y)|$ as the field itself.
Expand a domain-centered mixed correlation as
\[
 \sum_{T\subseteq\{1,\ldots,m\}}
   (-1)^{m-|T|}
   \prod_{k\notin T}m_D^b(y_k)
   \left\langle
     \prod_{i=1}^n\sigma(x_i)
     \prod_{k\in T}\epsilon(y_k)
   \right\rangle_D^b.
\]
Every term has the common conformal factor
$\prod_i|\phi'(x_i)|^{1/8}\prod_k|\phi'(y_k)|$.
Thus the actual-domain-centered nearest-neighbor comparison functions
have precisely the covariance stated in the theorem. The constants
$Z_\sigma^nZ_\epsilon^m$ do not depend on the domain and preserve this
identity. Homogeneous fixed and free boundary types are transported
unchanged. This completes every bulk correlation and covariance
assertion; no claim concerning coincident or boundary insertions, or
interface convergence, is required or used.
\end{proof}
